\originalchapter{来源覆盖与补充范围}
\originalsection{18.03 的课堂主线}
本册按数学依赖顺序重组, 不是逐讲逐页翻译. 官方目录列出36份课堂讲义, 第10、19、31次为考试位置, 目录无相应课堂讲义. 下表的数字是原课 session 编号, 不是本册页码. 所有原件逐文件页数、官方链接、SHA-256 和持有状态见 \texttt{source-manifest.json}.
\begin{longtable}{p{.16\textwidth}p{.47\textwidth}p{.25\textwidth}}
\toprule 原课讲次 & 数学内容 & 本册对应 \\\midrule\endhead
1 & 方向场、分离变量、存在唯一性 & 第1、2章 \\
2 & 数值方法 & 第1章误差衔接; 算法见数值分析册 \\
3--4 & 线性建模、积分因子 & 第2章 \\
5--7 & 复数、复指数、正弦响应 & 第4、5章; 复数代数作先修 \\
8--9 & 相线、线性与非线性 & 第3、9、10章 \\
11--12 & 特征根、阻尼振动 & 第4、5章 \\
13--16 & 指数响应、拍、共振与增益 & 第5章 \\
17--18 & LTI、RLC 与工程响应 & 第5、6章; 原工程实例不逐一重述 \\
20--22 & Fourier 级数、周期响应与共振 & 第7章 \\
23--25 & 阶跃、脉冲、卷积 & 第6章 \\
26--30 & Laplace、极点、传递函数 & 第6章; 差分补充见第13章 \\
32--34 & 线性系统、特征值、重根 & 第8章 \\
35--36 & 相图、模态、矩阵指数 & 第8、9章 \\
37--39 & 非线性、单摆、极限环与混沌入门 & 第9--11、13章 \\
\bottomrule
\end{longtable}
Miller 的 \textit{Supplementary Notes} 整本为155个 PDF 页面, 封面署名年份含2004、2006、2008和2010; 本次课程版本是2010. 核对 PDF 正文发现第1章与前言的分件含不同版本的文字, 因而另行保留这两份, 不按文件名强行去重. 其余25章与3个附录的数学内容对应整本, 单独的 OCW 末页和版式可能不同, 不把它们视为字节相同文件. 这28个重复分件登记官方链接并归并到整本持有. 官方资源网页的个别描述有误: 标记 chapter\_8 的原件实际为 \textit{RLC circuits}, 以 PDF 正文为准.

补充讲义第1--6节的记号、线性建模和复指数纳入第1、2、5章; 第7--15节的拍、RLC、响应、阻尼和 Wronskian 纳入第4、5章; 第16--24节的 Fourier、脉冲、卷积和极点纳入第6、7章; 第25--26节的系统与相图纳入第8、9章. 附录A的流行病主题在第3章给出 SIR 动力学与最终规模推导; 附录B的桥梁颤振历史、附录C的 phugoid 工程详例留作延伸, 未声称本册完整译出. 这三篇随整本归档, 其来源身份与教学覆盖范围分开记录.

Mattuck 的 \textit{Notes and Exercises} 以官方阅读目录署名为依据, 不将所有文件的 PDF 元数据中统一出现的 Miller 名字误作逐件作者. 笔记16件涉及图解、复数、积分形式、算子、输入响应、Laplace、线性系统 LS.1--6、图示系统及极限环. 其中 S 原件的实际标题为 Stability, 对应第5、8章, 不依据网页的含糊描述臆测其主题. 第12章的幂级数、正则奇点与边值内容为编者补充; 原练习目录虽含 power series 主题, 不把它冒称为已完整翻译的课堂讲义. 14份原练习及解答在调查清单中标为范围外链接, 不归档、不宣称逐题翻译. Jeremy Orloff 的差分与 Z 变换补充为14页, 第13章只采用线性响应及稳定极点衔接, 不扩展成离散控制课程.
\originalsection{12.006J 的动力系统补充}
16份文件按原课程组合讲次登记. 含出版社保留图片或经许可使用的第三方图片者, 整份原件仅保留官方链接; 本册不提取、描摹或重新标注这些图片. 所用相图由明确方程独立构造.
\begin{longtable}{p{.16\textwidth}p{.48\textwidth}p{.24\textwidth}}
\toprule 原课讲次 & 原内容 & 本册采用边界 \\\midrule\endhead
1 & 非线性与混沌引论 & 第3、13章概念衔接 \\
2--3 & 一维流与分岔 & 第3章 \\
4--5 & 单摆、二维稳定性 & 第9、10章 \\
6--7 & 耗散与体积收缩 & 第10、13章 \\
8--9 & 受迫振子、极限环 & 第5、11章 \\
10--11 & 二维分岔 & 第11章正规模型; 一般定理未全文证明 \\
12--14 & 动力系统谱分析 & 第7章仅线性周期响应; 随机谱留延伸 \\
15--16 & Poincar\'e 截面 & 第11章 \\
17--18 & 流体、Rayleigh--B\'enard 对流 & 留延伸; 偏微分方程不纳入 \\
19 & 奇异吸引子 & 第13章机制提示; 不作全局分类 \\
20--21 & Lorenz 方程 & 第13章平衡、吸收界及线性失稳 \\
22 & H\'enon 吸引子 & 留延伸; 未把二维映射全部理论纳入 \\
23 & 分形维数 & 留延伸; 未建立分形测度理论 \\
24 & Lyapunov 指数 & 第13章定义与变分增长; 不证明遍历极限存在 \\
25--27 & 倍周期通往混沌 & 第13章 logistic 二周期; 不证明通用常数 \\
28 & 间歇性、准周期 & 留延伸 \\
\bottomrule
\end{longtable}
\originalsection{编者补充与引用工具}
本册补全局部积分迭代、最大延续、连续依赖、矩阵指数、线性化稳定性的 Lyapunov 证明、紧集 LaSalle 论证、平面周期乘子与 Dulac 判据. 幂级数的正系数支配证明、边值 Green 核与相容条件、Lorenz 吸收界和三次稳定判据也作为补充论证明确写出. 例题及练习为本册整理或另设, 不复用商业教材题面.

有限维 Jordan 标准形、一致多项式逼近、基本积分定理、隐函数定理以及平面 Jordan 曲线和 Green 公式作为已声明的先修工具使用. Hartman--Grobman 与 Poincar\'e--Bendixson 的一般定理准确列条件后引用, 不声称已给完整拓扑证明. Van der Pol 极限环的全局唯一性使用核实的 LC 第5页 Levinson--Smith 引用版本, 同样明确为引用工具. 一般 Hopf 定理、稳定流形定理和混沌的全局存在证明不作为本册已证明结果. 数值绘图或抽样轨道不承担这些证明.
